% Part: first-order-logic 
% Chapter: axiomatic-deduction 
% Section: rules-and-proofs

\documentclass[../../../include/open-logic-section]{subfiles}

\begin{document}

\iftag{FOL}
      {\olfileid{fol}{axd}{rul}}
      {\olfileid{pl}{axd}{rul}}

\olsection{Aturan lan \usetoken{P}{derivation}}

\begin{explain}
  !!{derivation}s aksiomatik mbokmenawa dadi sistem !!{derivation}
  kang paling prasaja kanggo logika. !!^a{derivation} mung sawijining
  rerangken !!{formula}s. Supaya dianggep !!a{derivation}, saben
  !!{formula} ing rerangken kasebut kudu dadi conto sawijining
  aksioma, utawa kudu ndherek saka siji utawa luwih !!{formula}s kang
  ndhisiki ing rerangken miturut sawijining aturan inferensi.
  !!^a{derivation} !!{derive}s !!{formula} kang pungkasan.
\end{explain}

\begin{defn}[!!^{derivability}]
Yen $\Gamma$ iku himpunan !!{formula}s saka $\Lang L$, mula
\emph{!!{derivation}} saka~$\Gamma$ yaiku rerangken winates
$!A_1$, \dots,~$!A_n$ saka !!{formula}s, ing ngendi kanggo saben
$i \le n$ salah siji saka perkara iki lumaku:
\begin{enumerate}
\item $!A_i \in \Gamma$; utawa
\item $!A_i$ iku aksioma; utawa
\item $!A_i$ ndherek saka sawijining $!A_j$ (lan $!A_k$) kanthi
  $j < i$ (lan $k < i$) miturut aturan inferensi.
\end{enumerate}
\end{defn}

Apa kang dianggep minangka !!{derivation} kang bener gumantung marang
aturan inferensi endi kang diidini (lan mesthi wae apa kang dianggep
aksioma). Aturan inferensi iku pratelan yen-banjur kang ngandhani yen,
ing syarat tartamtu, sawijining langkah~$!A_i$ ing !!a{derivation}
iku langkah inferensi kang bener. Cathetan koreksi sumber OLPL-024:
sumber Inggris ngilangake tandha metavariabel formula saka langkah
kang lagi dirembug, sanadyan definisi sadurunge lan kabeh sebutan
sabanjure nggunakake tandha kasebut; tandha iku dibalekake ing kene.

\begin{defn}[Aturan inferensi]
\emph{Aturan inferensi} menehi syarat cukup kanggo apa kang dianggep
langkah inferensi kang bener ing !!a{derivation} saka~$\Gamma$.
\end{defn}

Upamane, amarga saben rerangken siji-unsur $!A$ kanthi $!A \in \Gamma$
kanthi langsung dianggep !!a{derivation}, pratelan iki bisa dadi aturan
inferensi kang prasaja banget:
\begin{quote}
  Yen $!A \in \Gamma$, mula $!A$ mesthi dadi langkah inferensi kang
  bener ing saben !!{derivation} saka~$\Gamma$.
\end{quote}
Semono uga, yen $!A$ iku salah siji aksioma, $!A$ dhewe dadi
!!a{derivation}, mula pratelan iki uga dadi aturan inferensi:
\begin{quote}
  Yen $!A$ iku aksioma, mula $!A$ dadi langkah inferensi kang bener.
\end{quote}
Bab iki dadi luwih narik kawigaten yen aturan inferensi nggunakake
!!{formula}s kang muncul sadurunge langkah kang ditimbang. Aturan iki
diarani \emph{modus ponens:}
\begin{quote}
  Yen $!B \lif !A$ lan $!B$ muncul luwih dhuwur ing !!{derivation},
  mula~$!A$ dadi langkah inferensi kang bener.
\end{quote}
Yen iki siji-sijine aturan inferensi, mula definisi !!{derivation} ing
ndhuwur padha karo iki: $!A_1$, \dots,~$!A_n$ iku !!a{derivation}
yen lan mung yen kanggo saben $i \le n$ salah siji saka perkara iki
lumaku:
\begin{enumerate}
\item $!A_i \in \Gamma$; utawa
\item $!A_i$ iku aksioma; utawa
\item kanggo sawijining $j < i$, $!A_j$ yaiku $!B \lif !A_i$, lan
  kanggo sawijining $k < i$, $!A_k$ yaiku~$!B$.
\end{enumerate}
Klausa pungkasan nyebut yen $!A_i$ ndherek saka~$!A_j$ ($!B \lif
!A_i$) lan $!A_k$ ($!B$) miturut modus ponens. Yen kita bisa mlaku
saka $1$ nganti~$n$, lan saben wektu kita nemokake !!a{formula}~$!A_i$
kang kalebu ing~$\Gamma$, dadi aksioma, utawa diarani langkah inferensi
kang bener dening aturan inferensi, mula kabeh rerangken dianggep
!!{derivation} kang bener.

\begin{defn}[!!^{derivability}]
!!{formula}~$!A$ \emph{bisa !!{derivable}} saka $\Gamma$, ditulis
$\Gamma \Proves !A$, yen ana !!a{derivation} saka $\Gamma$ kang
pungkasan ing $!A$.
\end{defn}

\begin{defn}[Teorema]
!!^a{formula}~$!A$ iku \emph{teorema} yen ana !!a{derivation}
saka himpunan kosong kang ngasilake~$!A$. Kita nulis $\Proves !A$
yen $!A$ iku teorema, lan $\Proves/ !A$ yen dudu teorema.
\end{defn}


\end{document}
